\section{The canonical delta-functor}
\label{section-canonical-delta-functor}

\noindent
The derived category should be the receptacle for the universal
cohomology functor. In order to state the result we use
the notion of a $\delta$-functor from an abelian category
into a triangulated category, see
Definition \ref{definition-delta-functor}.

\medskip\noindent
Consider the functor
$\text{Comp}(\mathcal{A}) \to K(\mathcal{A})$.
This functor is {\bf not} a $\delta$-functor in general.
The easiest way to see this is to consider a nonsplit
short exact sequence $0 \to A \to B \to C \to 0$
of objects of $\mathcal{A}$. Since
$\Hom_{K(\mathcal{A})}(C[0], A[1]) = 0$
we see that any distinguished triangle arising from
this short exact sequence would look like
$(A[0], B[0], C[0], a, b, 0)$. But the existence of such a
distinguished triangle in $K(\mathcal{A})$ implies
that the extension is split. A contradiction.

\medskip\noindent
It turns out that the functor
$\text{Comp}(\mathcal{A}) \to D(\mathcal{A})$ is a
$\delta$-functor. In order to see this we have to define
the morphisms $\delta$ associated to a short exact sequence
$$
0 \to A^\bullet \xrightarrow{a} B^\bullet \xrightarrow{b} C^\bullet \to 0
$$
of complexes in the abelian category $\mathcal{A}$.
Consider the cone $C(a)^\bullet$ of the morphism $a$.
We have $C(a)^n = B^n \oplus A^{n + 1}$ and we define
$q^n : C(a)^n \to C^n$ via the projection to $B^n$ followed
by $b^n$. Hence a morphism of complexes
$$
q : C(a)^\bullet \longrightarrow C^\bullet.
$$
It is clear that $q \circ i = b$ where $i$ is as in
Definition \ref{definition-cone}.
Note that, as $a^\bullet$ is injective in each degree,
the kernel of $q$ is identified with the cone of
$\text{id}_{A^\bullet}$ which is acyclic. Hence we see that
$q$ is a quasi-isomorphism. According to
Lemma \ref{lemma-the-same-up-to-isomorphisms}
the triangle
$$
(A, B, C(a), a, i, -p)
$$
is a distinguished triangle in $K(\mathcal{A})$.
As the localization functor
$K(\mathcal{A}) \to D(\mathcal{A})$ is
exact we see that $(A, B, C(a), a, i, -p)$ is a distinguished
triangle in $D(\mathcal{A})$. Since $q$ is a quasi-isomorphism
we see that $q$ is an isomorphism in $D(\mathcal{A})$.
Hence we deduce that
$$
(A, B, C, a, b, -p \circ q^{-1})
$$
is a distinguished triangle of $D(\mathcal{A})$.
This suggests the following lemma.

\begin{lemma}
\label{lemma-derived-canonical-delta-functor}
Let $\mathcal{A}$ be an abelian category. The functor
$\text{Comp}(\mathcal{A}) \to D(\mathcal{A})$
defined has the natural structure of a $\delta$-functor,
with
$$
\delta_{A^\bullet \to B^\bullet \to C^\bullet} = - p \circ q^{-1}
$$
with $p$ and $q$ as explained above. The same construction turns the
functors
$\text{Comp}^{+}(\mathcal{A}) \to D^{+}(\mathcal{A})$,
$\text{Comp}^{-}(\mathcal{A}) \to D^{-}(\mathcal{A})$, and
$\text{Comp}^b(\mathcal{A}) \to D^b(\mathcal{A})$
into $\delta$-functors.
\end{lemma}

\begin{proof}
We have already seen that this choice leads to a distinguished
triangle whenever given a short exact sequence of complexes.
We have to show that given a commutative diagram
$$
\xymatrix{
0 \ar[r] &
A^\bullet \ar[r]_a \ar[d]_f &
B^\bullet \ar[r]_b \ar[d]_g &
C^\bullet \ar[r] \ar[d]_h &
0 \\
0 \ar[r] &
(A')^\bullet \ar[r]^{a'} &
(B')^\bullet \ar[r]^{b'} &
(C')^\bullet \ar[r] &
0
}
$$
we get the desired commutative diagram of
Definition \ref{definition-delta-functor} (2).
By Lemma \ref{lemma-functorial-cone}
the pair $(f, g)$ induces a canonical morphism
$c : C(a)^\bullet \to C(a')^\bullet$. It is a simple computation
to show that $q' \circ c = h \circ q$ and
$f[1] \circ p = p' \circ c$. From this the result follows directly.
\end{proof}

\begin{lemma}
\label{lemma-derived-compare-triangles-ses}
Let $\mathcal{A}$ be an abelian category.
Let
$$
\xymatrix{
0 \ar[r] &
A^\bullet \ar[r] \ar[d] &
B^\bullet \ar[r] \ar[d] &
C^\bullet \ar[r] \ar[d] &
0 \\
0 \ar[r] &
D^\bullet \ar[r] &
E^\bullet \ar[r] &
F^\bullet \ar[r] &
0
}
$$
be a commutative diagram of morphisms of complexes
such that the rows are short exact sequences of complexes, and
the vertical arrows are quasi-isomorphisms.
The $\delta$-functor of
Lemma \ref{lemma-derived-canonical-delta-functor}
above maps the short exact sequences
$0 \to A^\bullet \to B^\bullet \to C^\bullet \to 0$
and
$0 \to D^\bullet \to E^\bullet \to F^\bullet \to 0$
to isomorphic distinguished triangles.
\end{lemma}

\begin{proof}
Trivial from the fact that $K(\mathcal{A}) \to D(\mathcal{A})$
transforms quasi-isomorphisms into isomorphisms and that the
associated distinguished triangles are functorial.
\end{proof}

\begin{lemma}
\label{lemma-derived-compare-triangles-split-case}
Let $\mathcal{A}$ be an abelian category. Let
$$
\xymatrix{
0 \ar[r] &
A^\bullet \ar[r] &
B^\bullet \ar[r] &
C^\bullet \ar[r] &
0
}
$$
be a short exact sequences of complexes.
Assume this short exact sequence is termwise split. Let
$(A^\bullet, B^\bullet, C^\bullet, \alpha, \beta, \delta)$
be the distinguished triangle of $K(\mathcal{A})$
associated to the sequence. The $\delta$-functor of
Lemma \ref{lemma-derived-canonical-delta-functor}
above maps the short exact sequences
$0 \to A^\bullet \to B^\bullet \to C^\bullet \to 0$
to a triangle isomorphic to the distinguished triangle
$$
(A^\bullet, B^\bullet, C^\bullet, \alpha, \beta, \delta).
$$
\end{lemma}

\begin{proof}
Follows from
Lemma \ref{lemma-the-same-up-to-isomorphisms}.
\end{proof}

\begin{remark}
\label{remark-truncation-distinguished-triangle}
Let $\mathcal{A}$ be an abelian category. Let $K^\bullet$ be a complex
of $\mathcal{A}$. Let $a \in \mathbf{Z}$. We claim there is a canonical
distinguished triangle
$$
\tau_{\leq a}K^\bullet \to K^\bullet \to \tau_{\geq a + 1}K^\bullet \to
(\tau_{\leq a}K^\bullet)[1]
$$
in $D(\mathcal{A})$. Here we have used the canonical truncation functors $\tau$
from Homology, Section \ref{homology-section-truncations}.
Namely, we first take the distinguished
triangle associated by our $\delta$-functor
(Lemma \ref{lemma-derived-canonical-delta-functor})
to the short exact sequence of complexes
$$
0 \to \tau_{\leq a}K^\bullet \to K^\bullet \to
K^\bullet/\tau_{\leq a}K^\bullet \to 0
$$
Next, we use that the map $K^\bullet \to \tau_{\geq a + 1}K^\bullet$
factors through a quasi-isomorphism
$K^\bullet/\tau_{\leq a}K^\bullet \to \tau_{\geq a + 1}K^\bullet$
by the description of cohomology groups in
Homology, Section \ref{homology-section-truncations}.
In a similar way we obtain canonical distinguished triangles
$$
\tau_{\leq a}K^\bullet \to \tau_{\leq a + 1}K^\bullet \to
H^{a + 1}(K^\bullet)[-a-1] \to (\tau_{\leq a}K^\bullet)[1]
$$
and
$$
H^a(K^\bullet)[-a] \to \tau_{\geq a}K^\bullet \to \tau_{\geq a + 1}K^\bullet
\to H^a(K^\bullet)[-a + 1]
$$
\end{remark}

\begin{lemma}
\label{lemma-trick-vanishing-composition}
Let $\mathcal{A}$ be an abelian category. Let
$$
K_0^\bullet \to K_1^\bullet \to \ldots \to K_n^\bullet
$$
be maps of complexes such that
\begin{enumerate}
\item $H^i(K_0^\bullet) = 0$ for $i > 0$,
\item $H^{-j}(K_j^\bullet) \to H^{-j}(K_{j + 1}^\bullet)$ is zero.
\end{enumerate}
Then the composition $K_0^\bullet \to K_n^\bullet$ factors through
$\tau_{\leq -n}K_n^\bullet \to K_n^\bullet$ in $D(\mathcal{A})$.
Dually, given maps of complexes
$$
K_n^\bullet \to K_{n - 1}^\bullet \to \ldots \to K_0^\bullet
$$
such that
\begin{enumerate}
\item $H^i(K_0^\bullet) = 0$ for $i < 0$,
\item $H^j(K_{j + 1}^\bullet) \to H^j(K_j^\bullet)$ is zero,
\end{enumerate}
then the composition $K_n^\bullet \to K_0^\bullet$ factors through
$K_n^\bullet \to \tau_{\geq n}K_n^\bullet$ in $D(\mathcal{A})$.
\end{lemma}

\begin{proof}
The case $n = 1$. Since $\tau_{\leq 0}K_0^\bullet = K_0^\bullet$
in $D(\mathcal{A})$ we can replace
$K_0^\bullet$ by $\tau_{\leq 0}K_0^\bullet$ and
$K_1^\bullet$ by $\tau_{\leq 0}K_1^\bullet$.
Consider the distinguished triangle
$$
\tau_{\leq -1}K_1^\bullet \to K_1^\bullet \to
H^0(K_1^\bullet)[0] \to (\tau_{\leq -1}K_1^\bullet)[1]
$$
(Remark \ref{remark-truncation-distinguished-triangle}).
The composition $K_0^\bullet \to K_1^\bullet \to H^0(K_1^\bullet)[0]$
is zero as it is equal to $K_0^\bullet \to H^0(K_0^\bullet)[0] \to
H^0(K_1^\bullet)[0]$ which is zero by assumption.
The fact that $\Hom_{D(\mathcal{A})}(K_0^\bullet, -)$
is a homological functor (Lemma \ref{lemma-representable-homological}),
allows us to find the desired factorization.
For $n = 2$ we get a factorization
$K_0^\bullet \to \tau_{\leq -1}K_1^\bullet$ by the case $n = 1$
and we can apply the case $n = 1$ to the map of complexes
$\tau_{\leq -1}K_1^\bullet \to \tau_{\leq -1}K_2^\bullet$
to get a factorization
$\tau_{\leq -1}K_1^\bullet \to \tau_{\leq -2}K_2^\bullet$.
The general case is proved in exactly the same manner.
\end{proof}




